% GENERATED FILE. Edit canonical lessons, then run npm run generate:books.
% canonical-source-hash: fnv1a64:910c5ae200181ae1
% canonical-lessons: HI-C372-daya, HI-C372-karuna, HI-C372-vafadar, HI-C372-pap, HI-C372-thamna

\chapter{Mercy and Loyalty}
\label{ch:hi-mercy-and-loyalty}
\hlchaptermodality{372}

\begin{chapteropening}
\textbf{By the end of this chapter:} \emph{I can talk about mercy, loyalty and wrongdoing.}

Complete the last lesson of chapter 372: I can talk about mercy, loyalty and wrongdoing.
\end{chapteropening}

\section[\texorpdfstring{\hi{दया} (dayā)}{दया (dayā)}]{\hi{दया} (dayā) — mercy, pity}

\label{lesson:HI-C372-daya}



\begin{quote}
Before the new one: say the Hindi for apart from, then the Hindi for to dig one's heels in.
\end{quote}

\subsection*{You'll want to know: \hi{दया}}
\textbf{\hi{दया}} — \emph{dayā} — \textquotedblleft{}mercy, pity\textquotedblright{}.

\begin{quote}
\hi{उसे} \hi{ग़रीबों} \hi{पर} \hi{दया} \hi{आती} \hi{है।}
\end{quote}

\begin{quote}
\emph{use ġarīboṁ par dayā ātī hai.}
\end{quote}

\begin{quote}
He feels pity for the poor.
\end{quote}

Feminine.

\begin{cousinweb}
\textbf{\hi{दया}} (\emph{dayā}) is a tatsama, borrowed from Sanskrit \emph{dayā}, \textquotedblleft{}sympathy, compassion\textquotedblright{} (Monier-Williams).
\end{cousinweb}

\begin{grammarlens}[title={What you've built}]
One more word for giving reasons and opinions.
\end{grammarlens}

\subsection*{Your turn}
\begin{itemize}
\raggedright
  \item \emph{Say it:} \emph{dayā}
  \item \emph{Say it:} \emph{dayā}, once more
  \item \emph{Say it:} \emph{use ġarīboṁ par dayā ātī hai.} — He feels pity for the poor.
  \item \emph{Recall:} say the Hindi for apart from, then the Hindi for to dig one's heels in, then say \emph{dayā} again
\end{itemize}

\begin{tcolorbox}[breakable,colback=teal!4,colframe=teal!35!black,title={Before you move on}]
What does \hi{दया} mean? (\textquotedblleft{}Mercy, pity\textquotedblright{}.) Say it once more.
\end{tcolorbox}
\section[\texorpdfstring{\hi{करुणा} (karuṇā)}{करुणा (karuṇā)}]{\hi{करुणा} (karuṇā) — compassion}

\label{lesson:HI-C372-karuna}



\begin{quote}
Before the new one: say the Hindi for to dig one's heels in, then the Hindi for mercy.
\end{quote}

\subsection*{You'll want to know: \hi{करुणा}}
\textbf{\hi{करुणा}} — \emph{karuṇā} — \textquotedblleft{}compassion\textquotedblright{}.

\begin{quote}
\hi{उसकी} \hi{आँखों} \hi{में} \hi{करुणा} \hi{थी।}
\end{quote}

\begin{quote}
\emph{uskī āṅkhoṁ meṁ karuṇā thī.}
\end{quote}

\begin{quote}
There was compassion in his eyes.
\end{quote}

Feminine.

\begin{cousinweb}
\textbf{\hi{करुणा}} (\emph{karuṇā}) is a tatsama, borrowed from Sanskrit \emph{karuṇā}, \textquotedblleft{}pity, compassion\textquotedblright{} (Monier-Williams).
\end{cousinweb}

\begin{grammarlens}[title={What you've built}]
One more word for giving reasons and opinions.
\end{grammarlens}

\subsection*{Your turn}
\begin{itemize}
\raggedright
  \item \emph{Say it:} \emph{karuṇā}
  \item \emph{Say it:} \emph{karuṇā}, once more
  \item \emph{Say it:} \emph{uskī āṅkhoṁ meṁ karuṇā thī.} — There was compassion in his eyes.
  \item \emph{Recall:} say the Hindi for to dig one's heels in, then the Hindi for mercy, then say \emph{karuṇā} again
\end{itemize}

\begin{tcolorbox}[breakable,colback=teal!4,colframe=teal!35!black,title={Before you move on}]
What does \hi{करुणा} mean? (\textquotedblleft{}Compassion\textquotedblright{}.) Say it once more.
\end{tcolorbox}
\section[\texorpdfstring{\hi{वफ़ादार} (vafādār)}{वफ़ादार (vafādār)}]{\hi{वफ़ादार} (vafādār) — loyal, faithful}

\label{lesson:HI-C372-vafadar}



\begin{quote}
Before the new one: say the Hindi for mercy, then the Hindi for compassion.
\end{quote}

\subsection*{You'll want to know: \hi{वफ़ादार}}
\textbf{\hi{वफ़ादार}} — \emph{vafādār} — \textquotedblleft{}loyal, faithful\textquotedblright{}.

\begin{quote}
\hi{कुत्ता} \hi{अपने} \hi{मालिक} \hi{का} \hi{वफ़ादार} \hi{होता} \hi{है।}
\end{quote}

\begin{quote}
\emph{kuttā apne mālik kā vafādār hotā hai.}
\end{quote}

\begin{quote}
A dog is loyal to its master.
\end{quote}

\begin{cousinweb}
\textbf{\hi{वफ़ादार}} (\emph{vafādār}) is usually traced to Persian \emph{vafādār}, \textquotedblleft{}keeping faith\textquotedblright{}. Its first part is Arabic \emph{wafāʾ}, \textquotedblleft{}faithfulness\textquotedblright{}.
\end{cousinweb}

\begin{grammarlens}[title={What you've built}]
One more word for giving reasons and opinions.
\end{grammarlens}

\subsection*{Your turn}
\begin{itemize}
\raggedright
  \item \emph{Say it:} \emph{vafādār}
  \item \emph{Say it:} \emph{vafādār}, once more
  \item \emph{Say it:} \emph{kuttā apne mālik kā vafādār hotā hai.} — A dog is loyal to its master.
  \item \emph{Recall:} say the Hindi for mercy, then the Hindi for compassion, then say \emph{vafādār} again
\end{itemize}

\begin{tcolorbox}[breakable,colback=teal!4,colframe=teal!35!black,title={Before you move on}]
What does \hi{वफ़ादार} mean? (\textquotedblleft{}Loyal, faithful\textquotedblright{}.) Say it once more.
\end{tcolorbox}
\section[\texorpdfstring{\hi{पाप} (pāp)}{पाप (pāp)}]{\hi{पाप} (pāp) — a sin, a wrong}

\label{lesson:HI-C372-pap}



\begin{quote}
Before the new one: say the Hindi for compassion, then the Hindi for loyal.
\end{quote}

\subsection*{You'll want to know: \hi{पाप}}
\textbf{\hi{पाप}} — \emph{pāp} — \textquotedblleft{}a sin, a wrong\textquotedblright{}.

\begin{quote}
\hi{झूठ} \hi{बोलना} \hi{पाप} \hi{है।}
\end{quote}

\begin{quote}
\emph{jhūṭh bolnā pāp hai.}
\end{quote}

\begin{quote}
Lying is a sin.
\end{quote}

Masculine.

\begin{cousinweb}
\textbf{\hi{पाप}} (\emph{pāp}) is a tatsama, borrowed from Sanskrit \emph{pāpa}, \textquotedblleft{}evil, wicked; a sin\textquotedblright{} (Monier-Williams).
\end{cousinweb}

\begin{grammarlens}[title={What you've built}]
One more word for giving reasons and opinions.
\end{grammarlens}

\subsection*{Your turn}
\begin{itemize}
\raggedright
  \item \emph{Say it:} \emph{pāp}
  \item \emph{Say it:} \emph{pāp}, once more
  \item \emph{Say it:} \emph{jhūṭh bolnā pāp hai.} — Lying is a sin.
  \item \emph{Recall:} say the Hindi for compassion, then the Hindi for loyal, then say \emph{pāp} again
\end{itemize}

\begin{tcolorbox}[breakable,colback=teal!4,colframe=teal!35!black,title={Before you move on}]
What does \hi{पाप} mean? (\textquotedblleft{}A sin, a wrong\textquotedblright{}.) Say it once more.
\end{tcolorbox}
\section[\texorpdfstring{\hi{थामना} (thāmnā)}{थामना (thāmnā)}]{\hi{थामना} (thāmnā) — to hold on to, to hold up}

\label{lesson:HI-C372-thamna}



\begin{quote}
Before the new one: say the Hindi for loyal, then the Hindi for a sin.
\end{quote}

\subsection*{You'll want to know: \hi{थामना}}
\textbf{\hi{थामना}} — \emph{thāmnā} — \textquotedblleft{}to hold on to, to hold up\textquotedblright{}.

\begin{quote}
\hi{बच्चे} \hi{ने} \hi{माँ} \hi{का} \hi{हाथ} \hi{थामा।}
\end{quote}

\begin{quote}
\emph{bacce ne māṁ kā hāth thāmā.}
\end{quote}

\begin{quote}
The child held on to his mother's hand.
\end{quote}

Past, with \hi{ने} (he or she …): \hi{उसने} \hi{मेरा} \hi{हाथ} \hi{थामा}, \emph{usne merā hāth thāmā}.

\begin{cousinweb}
\textbf{\hi{थामना}} (\emph{thāmnā}) is a tadbhava, inherited from Sanskrit \emph{stambhate}, \textquotedblleft{}props, stops\textquotedblright{} (CDIAL 13683).
\end{cousinweb}

\begin{grammarlens}[title={What you've built}]
One more word for giving reasons and opinions.
\end{grammarlens}

\subsection*{Your turn}
\begin{itemize}
\raggedright
  \item \emph{Say it:} \emph{thāmnā}
  \item \emph{Say it:} \emph{thāmnā}, once more
  \item \emph{Say it:} \emph{bacce ne māṁ kā hāth thāmā.} — The child held on to his mother's hand.
  \item \emph{Recall:} say the Hindi for loyal, then the Hindi for a sin, then say \emph{thāmnā} again
\end{itemize}

\begin{tcolorbox}[breakable,colback=teal!4,colframe=teal!35!black,title={Before you move on}]
What does \hi{थामना} mean? (\textquotedblleft{}To hold on to, to hold up\textquotedblright{}.) Say it once more.
\end{tcolorbox}
